


% Table generated by Excel2LaTeX from sheet 'Sheet1'
\begin{table*}[t]
	\centering
	%	\vspace{-8pt}
	\caption{A coarse summary of IL methods demonstrating different expert data modalities they can handle (learning from \underline{\textit{demonstrations}} or \underline{\textit{observations}}), disparate task settings they consider (learning from \underline{\textit{online}} environment interactions or pre-collected \underline{\textit{offline}} static dataset), the specific \underline{\textit{cross domain}} setting they assume (the transition dynamics between the learning environment and that of the expert behaviors are different).
    % , and the unique \underline{\textit{one-shot}} merit they desire (the learned policy can enable one-shot transfer to new imitation tasks). 
    We highlight that our contextual IL paradigm (CEIL) can naturally be applied to all the above IL settings without introducing any additional models or losses. 
	}
		\resizebox{\linewidth}{!}{
		\begin{tabular}{lccccccc}
			\toprule
			& \multicolumn{2}{c}{Expert Data $\{\pi_E(\tbtau)\}$} & \multicolumn{2}{c}{Task Setting} & \multirow{2}{*}{\makecell[c]{Cross Domain\\(dynamics)}} \\
			\cmidrule(lr){2-3} \cmidrule(lr){4-5} 
			& \multicolumn{1}{c}{Demonstrations} & \multicolumn{1}{c}{Observations} & \multicolumn{1}{c}{Online} & \multicolumn{1}{c}{Offline} \\% & \multicolumn{1}{c}{Cross Domain} & Zero-shot  \\
			\midrule
			Behavior Cloning (BC)  
			& {\scriptsize \CheckmarkBold} & {\scriptsize \textcolor{minetable1colorx}{\XSolidBrush}} & {\scriptsize \CheckmarkBold} & {\scriptsize \CheckmarkBold} & {\scriptsize \textcolor{minetable1colorx}{\XSolidBrush}}  \\
			GAIL~\citep{ho2016generative}
			& {\scriptsize \CheckmarkBold} & {\scriptsize \textcolor{minetable1colorx}{\XSolidBrush}} & {\scriptsize \CheckmarkBold} & {\scriptsize \textcolor{minetable1colorx}{\XSolidBrush}} & {\scriptsize \textcolor{minetable1colorx}{\XSolidBrush}}   \\
			% f-MAX  
			% & {\scriptsize \CheckmarkBold} & {\scriptsize \textcolor{minetable1colorx}{\XSolidBrush}} & {\scriptsize \CheckmarkBold} & {\scriptsize \textcolor{minetable1colorx}{\XSolidBrush}} & {\scriptsize \textcolor{minetable1colorx}{\XSolidBrush}}   \\
			AIRL~\citep{fu2017learning}  
			& {\scriptsize \CheckmarkBold} & {\scriptsize \textcolor{minetable1colorx}{\XSolidBrush}} & {\scriptsize \CheckmarkBold} & {\scriptsize \textcolor{minetable1colorx}{\XSolidBrush}} & {\scriptsize \textcolor{minetable1colorx}{\XSolidBrush}}   \\
			GAIfO~\citep{torabi2018generative} 
			& {\scriptsize \CheckmarkBold} & {\scriptsize \CheckmarkBold} & {\scriptsize \CheckmarkBold} & {\scriptsize \textcolor{minetable1colorx}{\XSolidBrush}} & {\scriptsize \textcolor{minetable1colorx}{\XSolidBrush}}   \\
			OPOLO~\citep{zhu2020off} 
			& {\scriptsize \CheckmarkBold} & {\scriptsize \CheckmarkBold} & {\scriptsize \CheckmarkBold} & {\scriptsize \textcolor{minetable1colorx}{\XSolidBrush}} & {\scriptsize \textcolor{minetable1colorx}{\XSolidBrush}}   \\
			IQ-Learn~\citep{garg2021iq}
			& {\scriptsize \CheckmarkBold} & {\scriptsize \CheckmarkBold} & {\scriptsize \CheckmarkBold} & {\scriptsize \CheckmarkBold} & {\scriptsize \textcolor{minetable1colorx}{\XSolidBrush}}  \\
			SQIL~\citep{reddy2019sqil}
			& {\scriptsize \CheckmarkBold} & {\scriptsize \textcolor{minetable1colorx}{\XSolidBrush}} & {\scriptsize \CheckmarkBold} & {\scriptsize \textcolor{minetable1colorx}{\XSolidBrush}} & {\scriptsize \textcolor{minetable1colorx}{\XSolidBrush}}  \\
			SAIL~\citep{liu2019state}
			& {\scriptsize \CheckmarkBold} & {\scriptsize \CheckmarkBold} & {\scriptsize \CheckmarkBold} & {\scriptsize \textcolor{minetable1colorx}{\XSolidBrush}} & {\scriptsize \CheckmarkBold}  \\
			ValueDICE~\citep{kostrikov2019imitation}
			& {\scriptsize \CheckmarkBold} & {\scriptsize \textcolor{minetable1colorx}{\XSolidBrush}} & {\scriptsize \CheckmarkBold} & {\scriptsize \textcolor{minetable1colorx}{\XSolidBrush}} & {\scriptsize \textcolor{minetable1colorx}{\XSolidBrush}}  \\
			DemoDICE~\citep{kim2022demodice}
			& {\scriptsize \CheckmarkBold} & {\scriptsize \textcolor{minetable1colorx}{\XSolidBrush}} & {\scriptsize \textcolor{minetable1colorx}{\XSolidBrush}} & {\scriptsize \CheckmarkBold} & {\scriptsize \textcolor{minetable1colorx}{\XSolidBrush}} \\
			ORIL~\citep{zolna2020offline}
			& {\scriptsize \CheckmarkBold} & {\scriptsize \CheckmarkBold} & {\scriptsize \textcolor{minetable1colorx}{\XSolidBrush}} & {\scriptsize \CheckmarkBold} & {\scriptsize \textcolor{minetable1colorx}{\XSolidBrush}}  \\
			SMODICE~\citep{ma2022smodice}
			& {\scriptsize \CheckmarkBold} & {\scriptsize \CheckmarkBold} & {\scriptsize \textcolor{minetable1colorx}{\XSolidBrush}} & {\scriptsize \CheckmarkBold} & {\scriptsize \CheckmarkBold} \\
			\toprule
			CEIL (ours)  
			& {\scriptsize \CheckmarkBold} & {\scriptsize \CheckmarkBold} & {\scriptsize \CheckmarkBold} & {\scriptsize \CheckmarkBold} & {\scriptsize \CheckmarkBold}   \\
			\bottomrule
		\end{tabular}%
			}
	% \label{tab:ceil-on-various-IL-settings}%
	%	\vspace{-9pt}
\end{table*}%% {\scriptsize \textcolor{minetable1colorx}{\XSolidBrush}}    & {\scriptsize \CheckmarkBold  }





Imitation learning (IL) methods can be categorized into two groups: learning from demonstrations (LfD) and learning from observations (LfO). LfD methods work with expert demonstrations (given with expert actions) which are typically collected by teleoperation or kinesthetic teaching. LfO methods learn from expert observations collected by observing the expert performing the task. 

\subsection{Online Learning from Demonstrations (LfD)}

LfD methods learn from expert demonstrations, which are typically collected by teleoperation or kinesthetic teaching. 
Behavior cloning (BC)~\cite{pomerleau1989alvinn,bojarski2016end} is a simple and effective LfD method that learns a policy by maximizing the probability of state-action pairs in the demonstrations. However, BC is prone to compounding errors and cannot recover from mistakes. 
Inverse reinforcement learning (IRL)~\cite{ng2000algorithms,abbeel2004apprenticeship} is another popular LfD method that learns a reward function from expert demonstrations and then learns a policy that maximizes the expected reward. However, IRL methods are typically computationally expensive and require access to the expert's policy. 
Generative adversarial imitation learning (GAIL)~\cite{ho2016generative} is a popular LfD method that learns a policy based on the feedback from a well-trained discriminator to act as similar as possible to an optimal policy. However, GAIL is prone to compounding errors and cannot recover from mistakes. 
Based on GAIL, AIRL~\citep{fu2017learning} modifies the form of the discriminator to explicitly include a parameterized reward function, with which the trained policy is more generalize and less sensitive to dynamics. More specifically, AIRL derives a state-dependent reward function enables off-dynamics imitation learning from cross-domain demonstrations.
SQIL~\citep{reddy2019sqil} is a simple alternative that gives a reward of one to the expert transitions and a reward of zero to the transitions stored in buffer. Essentially, SQIL incentivizes the agent to match demonstrations over a long horizon and return to demonstrated states.
ValueDICE~\citep{kostrikov2019imitation} is an off-policy imitation learning algorithm that transforms the traditional distribution ratio estimation objective to allow for completely off-policy data, providing data-efficiency and eliminating the need for a separate RL optimization.

\subsection{Online Learning from Observation}  
Learning from observation (LfO) enables learning directly from state-only demonstrations without access to the demonstrator's actions. 
GAIfO~\citep{torabi2018generative} modifies the discriminator of GAIL to only judge whether the generated trajectory matches the expert trajectory in terms of states, allowing GAIfO to accomplish LfO tasks in a GAN-based framework.
SAIL~\citep{liu2019state} trains the imitator to follow state sequences in expert demonstrations, incorporating both local and global perspectives into a reinforcement learning framework through a regularized policy update objective, and has been shown to outperform current imitation learning methods in settings where the expert and imitator have different dynamics models.
I2L~\citep{gangwani2020state} addresses the issue of expert demonstrations being unavailable or having discrepancies with the imitator MDP by introducing an indirection step and solving two subproblems iteratively through divergence minimization, and is particularly effective when there is a transition dynamics mismatch between the expert and imitator MDPs.
OPOLO~\citep{zhu2020off} proposes a sample-efficient LfO approach that enables off-policy optimization in a principled manner, using an inverse action model to assist distribution matching from the perspective of mode-covering, and achieves SOTA results in terms of sample-efficiency and asymptotic performance on challenging locomotion tasks.
GARAT~\citep{desai2020imitation} addresses the problem of transferring policies learned in a source environment to a target with different dynamics by re-purposing state-of-the-art approaches from the LfO literature to develop a new algorithm based on generative adversarial reinforced action transformation, which achieves higher returns in the target environment compared to existing black-box transfer methods in several domains with mismatched dynamics.
\citet{raychaudhuri2021cross} proposes a framework for imitation learning and domain transfer that allows for effective learning from unpaired multi-domain demonstrations. The framework incorporates Markov decision processes for inference and achieves a comparable level of effectiveness in the produced transferred demonstrations relative to self-demonstrations.

\subsection{Offline Imitation Learning} 
DemoDICE~\citep{kim2022demodice} is an offline DICE-based imitation learning method utilizing both expert and imperfect demonstrations to match the state distribution of expert policy while penalizing deviation from overall demonstrations, leading to a stabilized learning process compared to recent adversarial training objective-based IL algorithms.
ORIL~\citep{zolna2020offline} discusses offline learning from demonstrations and unlabeled experience in the context of behavior cloning, offline reinforcement learning, and positive-unlabeled learning. ORIL learns purely from logged data, without online deployment during learning, and uses all available demonstration and unlabeled set data.
IQ-Learn~\citep{garg2021iq} learns a single Q-function that implicitly represents both reward and policy, avoiding adversarial training, and achieving state-of-the-art results in offline and online imitation learning settings.
SMODICE~\citep{ma2022smodice} enable offline imitation learning from observations and examples that utilizes state-occupancy matching and regularization to generalize prior DICE-based offline methods. The proposed method solves an optimization problem that depends on samples from observations and examples, and obtains optimal state-occupancy importance weights for the offline data. 


